%============================================================
% BAB 1: PENDAHULUAN
%   1.1 Latar Belakang
%   1.2 Rumusan Masalah
%   1.3 Batasan Masalah
%   1.4 Tujuan
%       1.4.1 Tujuan Proyek [Nama Proyek 1]
%       1.4.2 Tujuan Proyek [Nama Proyek 2] (hapus jika hanya 1 proyek)
%   1.5 Manfaat
%       1.5.1 Manfaat Proyek [Nama Proyek 1]
%       1.5.2 Manfaat Proyek [Nama Proyek 2] (hapus jika hanya 1 proyek)
%   1.6 Sistematika Penulisan
%============================================================
\chapter{PENDAHULUAN}

\section{LATAR BELAKANG}

Deskripsikan latar belakang pemilihan tempat dan topik kerja praktik. Jelaskan
konteks industri atau institusi tempat KP dilaksanakan, permasalahan yang ada,
relevansi dengan program studi, serta gambaran umum proyek yang dikerjakan.
Latar belakang yang baik mampu menjelaskan urgensi dan motivasi pelaksanaan
kerja praktik ini.

\section{RUMUSAN MASALAH}

Berdasarkan latar belakang yang telah dijelaskan, maka permasalahan yang akan
dibahas dalam laporan kerja praktik ini dapat dirumuskan sebagai berikut.

\begin{enumerate}[label=\alph*., leftmargin=*, itemsep=0pt]
  \item Rumusan masalah pertama.
  \item Rumusan masalah kedua.
  \item Rumusan masalah ketiga.
\end{enumerate}

\section{BATASAN MASALAH}

Untuk memfokuskan pembahasan pada laporan kerja praktik ini, maka ditetapkan
batasan masalah sebagai berikut:

\begin{enumerate}[label=\alph*., leftmargin=*, itemsep=0pt]
  \item Pembahasan difokuskan pada kegiatan yang dilakukan selama periode
        \PeriodeKP\ di \NamaPerusahaan.
  \item Batasan teknis kedua (misalnya: hanya mencakup sisi frontend/backend).
  \item Batasan ketiga (aspek yang tidak dibahas dalam laporan ini).
  \item Periode pengamatan dan keterlibatan dalam proyek dibatasi selama masa
        kerja praktik.
\end{enumerate}

\section{TUJUAN}

Bagian ini menjelaskan tujuan dari pelaksanaan proyek kerja praktik yang
dilakukan di \NamaPerusahaan. Tujuan yang dirumuskan tidak hanya menggambarkan
hasil yang ingin dicapai dari sisi teknis, tetapi juga mencerminkan kontribusi
sistem terhadap efisiensi kerja dan nilai praktis dalam dunia industri.

% Jika hanya satu proyek, hapus sub-section 1.4.2 dan ubah 1.4.1 jadi narasi langsung.
\subsection{Tujuan Proyek [Nama Proyek 1]}

\begin{itemize}[leftmargin=*, itemsep=0pt]
  \item Tujuan pertama proyek.
  \item Tujuan kedua proyek.
  \item Tujuan ketiga proyek.
\end{itemize}

\subsection{Tujuan Proyek [Nama Proyek 2]}

\begin{itemize}[leftmargin=*, itemsep=0pt]
  \item Tujuan pertama proyek kedua.
  \item Tujuan kedua proyek kedua.
\end{itemize}

\section{MANFAAT}

Pelaksanaan kerja praktik di \NamaPerusahaan\ tidak hanya bertujuan untuk
memenuhi kewajiban akademik, tetapi juga memberikan manfaat nyata baik dari
sisi teknis, operasional, maupun strategis. Berikut adalah manfaat utama dari
masing-masing proyek:

% Jika hanya satu proyek, hapus sub-section 1.5.2 dan ubah 1.5.1 jadi narasi langsung.
\subsection{Manfaat Proyek [Nama Proyek 1]}

\begin{itemize}[leftmargin=*, itemsep=0pt]
  \item \textbf{Manfaat Pertama.} Deskripsi manfaat pertama.
  \item \textbf{Manfaat Kedua.} Deskripsi manfaat kedua.
  \item \textbf{Manfaat Ketiga.} Deskripsi manfaat ketiga.
\end{itemize}

\subsection{Manfaat Proyek [Nama Proyek 2]}

\begin{itemize}[leftmargin=*, itemsep=0pt]
  \item \textbf{Manfaat Pertama.} Deskripsi manfaat pertama proyek kedua.
  \item \textbf{Manfaat Kedua.} Deskripsi manfaat kedua proyek kedua.
\end{itemize}

\section{SISTEMATIKA PENULISAN}

Berikut sistematika penulisan pada laporan kerja praktik ini:

\input{\SistematikaFile}
